\documentclass[11pt]{article}

\usepackage[margin=1in]{geometry}
\usepackage{amsmath,amssymb,mathtools}
\usepackage{bm}
\usepackage{booktabs}
\usepackage{array}
\usepackage{microtype}
\usepackage[hidelinks]{hyperref}

\newcommand{\E}{\mathbb{E}}
\newcommand{\Var}{\operatorname{var}}
\newcommand{\Cov}{\operatorname{cov}}
\newcommand{\Normal}{\mathcal{N}}
\newcommand{\diag}{\operatorname{diag}}
\newcommand{\trans}{\intercal}
\newcommand{\Olayer}{\mathtt{O}}

\title{Event-Based Approximate Gaussian Categorical Inference\\
within the TAGI Framework}
\author{}
\date{}

\begin{document}
\maketitle

\section{Purpose and relationship with TAGI}

Tractable Approximate Gaussian Inference (TAGI) propagates Gaussian moments
through a neural network and computes the posterior moments of the output units
$\bm Z^{(\Olayer)}$ before recursively propagating those moments backward using
Gaussian conditional (RTS) equations.  For a continuous observation, the
standard TAGI observation model is
\begin{equation}
    \bm Y = \bm Z^{(\Olayer)} + \bm V,
    \qquad
    \bm V \sim \Normal(\bm 0,\bm\Sigma_{\bm V}).
\end{equation}
This model is not appropriate for a one-hot categorical observation.  Event
AGCI replaces only this output observation model.  The TAGI forward pass and
the recursive TAGI backward pass remain unchanged.

For a class observation $Y=c$, Event AGCI performs the following sequence:
\begin{equation}
\boxed{
\begin{gathered}
\text{TAGI prior for }\bm Z^{(\Olayer)}
\;\longrightarrow\;
\text{class-winning event}
\;\longrightarrow\;
\text{event-conditioned moments}
\\[-0.2em]
\longrightarrow\;
\text{TAGI RTS backward pass}
\end{gathered}}
\end{equation}

\section{TAGI prior for the output units}

For one input, the TAGI forward pass supplies
\begin{equation}
 f(\bm z^{(\Olayer)})
 =
 \Normal\!\left(
 \bm z^{(\Olayer)};
 \bm\mu_{\bm Z^{(\Olayer)}},
 \bm\Sigma_{\bm Z^{(\Olayer)}}
 \right).
 \label{eq:tagi-output-prior}
\end{equation}
For $K$ classes, $\bm Z^{(\Olayer)}\in\mathbb R^K$ contains one latent
utility per class.  The standard diagonal TAGI approximation is
\begin{equation}
 \bm\Sigma_{\bm Z^{(\Olayer)}}
 =
 \diag\!\left(
 (\sigma_{Z_1^{(\Olayer)}})^2,
 \ldots,
 (\sigma_{Z_K^{(\Olayer)}})^2
 \right).
 \label{eq:diagonal-output-prior}
\end{equation}
For concision, define
\begin{equation}
 \bm\mu \equiv \bm\mu_{\bm Z^{(\Olayer)}},
 \qquad
 \bm\Sigma \equiv \bm\Sigma_{\bm Z^{(\Olayer)}},
 \qquad
 v_k \equiv (\sigma_{Z_k^{(\Olayer)}})^2.
\end{equation}

\section{Initialization of a TAGI last layer over frozen features}
\label{sec:initialization}

Let a deterministic neural network be frozen and let its final feature vector
be
\begin{equation}
 \bm H=f_{\mathrm{det}}(\bm x)\in\mathbb R^D.
\end{equation}
Only the final affine layer is treated probabilistically:
\begin{equation}
 Z_k^{(\Olayer)}
 =
 \bm W_k^\trans\bm H+B_k,
 \qquad k=1,\ldots,K.
 \label{eq:frozen-last-layer}
\end{equation}
The diagonal Gaussian prior is
\begin{equation}
 q_0(\bm W,\bm B)
 =
 \prod_{d=1}^{D}\prod_{k=1}^{K}
 \Normal(W_{k,d};\mu_{W_{k,d},0},(\sigma_{W_{k,d},0})^2)
 \prod_{k=1}^{K}
 \Normal(B_k;\mu_{B_k,0},(\sigma_{B_k,0})^2).
 \label{eq:last-layer-prior}
\end{equation}

\subsection{Prior means}

For a newly initialized classifier with no prior preference among classes, use
\begin{equation}
 \boxed{
 \mu_{W_{k,d},0}=0,
 \qquad
 \mu_{B_k,0}=0.}
 \label{eq:zero-prior-means}
\end{equation}
This produces equal utility means and, under class-symmetric variances, the
uniform prior predictive distribution
\begin{equation}
 P(Y=k)=\frac{1}{K}.
\end{equation}
Random parameter means are not needed to break symmetry: the first labeled
observations create different class-specific posterior updates.  Random means
would instead center the Bayesian prior on one arbitrary deterministic
classifier.

The multinomial-probit likelihood is invariant to a common utility offset.  A
zero-mean initialization fixes the location of the prior mean, but a fully
zero-sum random parameterization would require cross-class covariance or a
$K-1$ contrast parameterization.  The standard diagonal family fixes only the
mean gauge.

\subsection{Feature preprocessing}

The scale of an isotropic weight prior is meaningful only relative to the
feature coordinates.  Estimate all feature statistics using training data
only.  A recommended preprocessing sequence is:
\begin{enumerate}
 \item estimate and subtract the training-feature mean;
 \item optionally whiten the features or apply one global root-mean-square
 scale;
 \item freeze this transform and reuse it for validation and test inputs;
 \item initialize the TAGI prior in the transformed coordinates.
\end{enumerate}
Define the mean per-coordinate feature energy
\begin{equation}
 e_H
 =
 \frac{1}{D}\E_{\bm H}\!\left[\lVert\bm H\rVert^2\right].
 \label{eq:feature-energy}
\end{equation}

\subsection{Prior variances from a dimensionless utility scale}

Use a class-symmetric isotropic weight prior
\begin{equation}
 (\sigma_{W_{k,d},0})^2
 =
 \frac{g_W^2}{D},
 \label{eq:weight-prior-scale}
\end{equation}
and a bias prior
\begin{equation}
 (\sigma_{B_k,0})^2=\sigma_{B,0}^2.
 \label{eq:bias-prior-scale}
\end{equation}
The expected prior utility variance is then
\begin{align}
 \E_{\bm H}\!\left[
 (\sigma_{Z_k^{(\Olayer)},0})^2
 \right]
 &=
 \E_{\bm H}\!\left[
 \sum_{d=1}^{D}H_d^2\frac{g_W^2}{D}
 +\sigma_{B,0}^2
 \right]
 \\
 &=
 g_W^2e_H+\sigma_{B,0}^2.
 \label{eq:expected-prior-utility-variance}
\end{align}
Define the dimensionless prior utility ratio
\begin{equation}
 \boxed{
 \kappa^2
 =
 \frac{
 \E_{\bm H}[(\sigma_{Z_k^{(\Olayer)},0})^2]
 }{\tau^2}.}
 \label{eq:kappa-definition}
\end{equation}
For fixed $\tau$, selected $\kappa$, and selected bias variance, the weight
scale is
\begin{equation}
 \boxed{
 g_W
 =
 \sqrt{
 \frac{
 \kappa^2\tau^2-\sigma_{B,0}^2
 }{e_H}
 },}
 \qquad
 \kappa^2\tau^2>\sigma_{B,0}^2.
 \label{eq:weight-gain-initialization}
\end{equation}
Thus $\tau$ fixes the likelihood unit, $\kappa$ specifies prior uncertainty
relative to that unit, and $e_H$ converts the dimensionless prior choice into
a parameter variance.  If an implementation scales the bias variance by fan
in, $(\sigma_{B_k,0})^2=g_B^2/D$, substitute $g_B^2/D$ for
$\sigma_{B,0}^2$ in Equations
\eqref{eq:expected-prior-utility-variance}--\eqref{eq:weight-gain-initialization}.

A small $\kappa$ represents a concentrated prior and produces small posterior
movement.  A large $\kappa$ represents greater prior uncertainty and permits
larger evidence-driven updates.  Choosing $\kappa$ using validation likelihood
is an empirical-Bayes hyperparameter selection and should be reported as such.

\subsection{Pretrained last-layer means}

If the deterministic classifier already contains fitted final-layer
parameters $(\bm W_\star,\bm B_\star)$, a warm-start prior may use
\begin{equation}
 \bm\mu_{\bm W,0}=\bm W_\star,
 \qquad
 \bm\mu_{\bm B,0}=\bm B_\star.
\end{equation}
Its covariance should quantify uncertainty around that fitted solution.  In
decreasing order of fidelity, use a covariance retained from the original
training procedure, a Laplace or Fisher approximation, or a declared diagonal
effective prior variance.  A trusted pretrained mean should not normally be
combined with the same diffuse covariance used for a new zero-mean layer.

If the pretrained classifier was fitted on the same observations subsequently
assimilated by Event AGCI, those observations are counted twice.  Such a warm
start is an empirical-Bayes or recycled-evidence approximation, not ordinary
sequential Bayes.  Before reuse, subtract the common class-wise mean from the
pretrained weights and biases; this removes the deterministic common utility
offset without changing class probabilities.

\section{Gaussian random-utility observation model}

Introduce an independent Gaussian decision-noise vector
\begin{equation}
 \bm E \sim \Normal(\bm 0,\tau^2\mathbf I_K),
 \qquad \tau>0,
 \label{eq:decision-noise}
\end{equation}
and define the noisy utilities
\begin{equation}
 \bm A = \bm Z^{(\Olayer)} + \bm E.
 \label{eq:noisy-utilities}
\end{equation}
It follows that
\begin{equation}
 f(\bm a)
 =
 \Normal(\bm a;\bm\mu_{\bm A},\bm\Sigma_{\bm A}),
 \qquad
 \bm\mu_{\bm A}=\bm\mu,
 \qquad
 \bm\Sigma_{\bm A}=\bm\Sigma+\tau^2\mathbf I_K.
 \label{eq:utility-prior}
\end{equation}
The observed categorical class is the index of the largest noisy utility:
\begin{equation}
 Y=c
 \quad\Longleftrightarrow\quad
 A_c>A_j \quad\text{for every }j\ne c.
 \label{eq:winning-event}
\end{equation}
The one-hot representation of the observation is $\bm y=\bm e_c$.  Equation
\eqref{eq:winning-event} defines a multinomial-probit categorical likelihood.
It guarantees non-negative class probabilities that sum to one without
interpreting the Gaussian output means as probabilities.

The scale $\tau$ is fixed and must be identical during training and prediction.
A convenient reference choice is $\tau=1$, which fixes the latent-utility unit.

\section{Probability of a class-winning event}

Under the diagonal prior in \eqref{eq:diagonal-output-prior}, the noisy
utilities are independent, with
\begin{equation}
 A_k\sim\Normal(\mu_k,r_k^2),
 \qquad
 r_k^2=v_k+\tau^2.
 \label{eq:utility-marginal}
\end{equation}
Conditioning temporarily on $A_c=a$, define
\begin{equation}
 \alpha_j(a)=\frac{a-\mu_j}{r_j},
 \qquad j\ne c.
\end{equation}
The probability that every competing utility is below $a$ is
\begin{equation}
 P(A_j<a,\ \forall j\ne c\mid A_c=a)
 =
 \prod_{j\ne c}\Phi\!\left(\alpha_j(a)\right),
\end{equation}
where $\Phi(\cdot)$ is the standard Gaussian cumulative distribution function.
Therefore, the predictive probability of class $c$ is the one-dimensional
integral
\begin{equation}
 p_c
 \equiv P(Y=c)
 =
 \int_{-\infty}^{\infty}
 \Normal(a;\mu_c,r_c^2)
 \prod_{j\ne c}\Phi\!\left(\alpha_j(a)\right)
 \,da.
 \label{eq:class-probability}
\end{equation}
Although the class event contains $K-1$ inequalities, independence of the
incoming utilities reduces its probability to one dimension.

Define the unnormalized event density
\begin{equation}
 w_c(a)
 =
 \Normal(a;\mu_c,r_c^2)
 \prod_{j\ne c}\Phi\!\left(\alpha_j(a)\right).
\end{equation}
The event-conditioned density of the winning noisy utility is
\begin{equation}
 q_c(a)
 \equiv f(a_c=a\mid Y=c)
 =\frac{w_c(a)}{p_c}.
 \label{eq:winner-density}
\end{equation}

\section{Event-conditioned moments of the noisy utilities}

\subsection{Winning utility}

The first two moments of the winning utility are
\begin{align}
 \mu_{A_c\mid\bm y}
 &=
 \E[A_c\mid Y=c]
 =
 \int a\,q_c(a)\,da,
 \label{eq:winner-mean}
 \\
 \E[A_c^2\mid Y=c]
 &=
 \int a^2 q_c(a)\,da.
 \label{eq:winner-second}
\end{align}
Hence
\begin{equation}
 (\sigma_{A_c\mid\bm y})^2
 =
 \E[A_c^2\mid Y=c]
 -\mu_{A_c\mid\bm y}^2.
 \label{eq:winner-variance}
\end{equation}

\subsection{Losing utilities}

For $j\ne c$, conditioning on $A_c=a$ and $Y=c$ implies $A_j<a$.
Consequently, $A_j$ follows an upper-truncated Gaussian.  Define the inverse
Mills ratio
\begin{equation}
 \lambda_j(a)
 =
 \frac{\phi(\alpha_j(a))}{\Phi(\alpha_j(a))},
 \label{eq:mills-ratio}
\end{equation}
where $\phi(\cdot)$ is the standard Gaussian density.  The conditional mean
and variance at a fixed winner value $a$ are
\begin{align}
 m_j(a)
 &\equiv
 \E[A_j\mid A_j<a]
 =
 \mu_j-r_j\lambda_j(a),
 \label{eq:loser-local-mean}
 \\
 s_j^2(a)
 &\equiv
 \Var(A_j\mid A_j<a)
 \\
 &=
 r_j^2\left[
 1-\alpha_j(a)\lambda_j(a)-\lambda_j(a)^2
 \right].
 \label{eq:loser-local-variance}
\end{align}
Integrating with respect to the event-conditioned winner density gives
\begin{align}
 \mu_{A_j\mid\bm y}
 &=
 \int m_j(a)q_c(a)\,da,
 \label{eq:loser-mean}
 \\
 \E[A_j^2\mid Y=c]
 &=
 \int\left[s_j^2(a)+m_j(a)^2\right]q_c(a)\,da,
 \label{eq:loser-second}
 \\
 (\sigma_{A_j\mid\bm y})^2
 &=
 \E[A_j^2\mid Y=c]-\mu_{A_j\mid\bm y}^2.
 \label{eq:loser-variance}
\end{align}

\subsection{Dense covariance created by the class event}

Although the noisy utilities are independent before observing the class, the
shared winning event creates cross-class covariance.  The required cross
moments are
\begin{align}
 \E[A_cA_j\mid Y=c]
 &=
 \int a\,m_j(a)q_c(a)\,da,
 &&j\ne c,
 \label{eq:winner-loser-cross}
 \\
 \E[A_jA_\ell\mid Y=c]
 &=
 \int m_j(a)m_\ell(a)q_c(a)\,da,
 &&j\ne\ell,\quad j,\ell\ne c.
 \label{eq:loser-loser-cross}
\end{align}
Collecting the first moments in $\bm\mu_{\bm A\mid\bm y}$ and the second
moments in $\E[\bm A\bm A^\trans\mid\bm y]$ yields
\begin{equation}
 \bm\Sigma_{\bm A\mid\bm y}
 =
 \E[\bm A\bm A^\trans\mid\bm y]
 -
 \bm\mu_{\bm A\mid\bm y}
 \bm\mu_{\bm A\mid\bm y}^\trans.
 \label{eq:event-utility-covariance}
\end{equation}
When the subsequent TAGI representation retains only marginal variances,
Equations \eqref{eq:winner-loser-cross}--\eqref{eq:loser-loser-cross} need not
be evaluated; only the diagonal of \eqref{eq:event-utility-covariance} is
required.

\section{Event-conditioned posterior moments of the TAGI output}

Before observing the class, $\bm Z^{(\Olayer)}$ and $\bm A$ are jointly
Gaussian.  Their cross-covariance is
\begin{equation}
 \bm\Sigma_{\bm Z^{(\Olayer)}\bm A}
 =
 \bm\Sigma_{\bm Z^{(\Olayer)}}.
\end{equation}
Define the Gaussian conditional gain
\begin{equation}
 \boxed{
 \mathbf J_{\bm Z^{(\Olayer)}}^{\bm A}
 =
 \bm\Sigma_{\bm Z^{(\Olayer)}\bm A}
 \bm\Sigma_{\bm A}^{-1}
 =
 \bm\Sigma_{\bm Z^{(\Olayer)}}
 \left(
 \bm\Sigma_{\bm Z^{(\Olayer)}}+\tau^2\mathbf I_K
 \right)^{-1}.}
 \label{eq:z-a-gain}
\end{equation}
Applying the law of total expectation to the Gaussian conditional
$f(\bm z^{(\Olayer)}\mid\bm a)$ gives
\begin{equation}
 \boxed{
 \bm\mu_{\bm Z^{(\Olayer)}\mid\bm y}
 =
 \bm\mu_{\bm Z^{(\Olayer)}}
 +
 \mathbf J_{\bm Z^{(\Olayer)}}^{\bm A}
 \left(
 \bm\mu_{\bm A\mid\bm y}-\bm\mu_{\bm A}
 \right).}
 \label{eq:event-output-mean}
\end{equation}
The law of total covariance gives
\begin{equation}
 \boxed{
 \bm\Sigma_{\bm Z^{(\Olayer)}\mid\bm y}
 =
 \bm\Sigma_{\bm Z^{(\Olayer)}}
 +
 \mathbf J_{\bm Z^{(\Olayer)}}^{\bm A}
 \left(
 \bm\Sigma_{\bm A\mid\bm y}-\bm\Sigma_{\bm A}
 \right)
 \left(\mathbf J_{\bm Z^{(\Olayer)}}^{\bm A}\right)^\trans.}
 \label{eq:event-output-covariance}
\end{equation}
Equations \eqref{eq:event-output-mean} and
\eqref{eq:event-output-covariance} have the same RTS form as the TAGI
backward equations.  The difference is that the posterior moments of
$\bm A\mid\bm y$ are produced by the class-winning event instead of a linear
Gaussian observation.

For diagonal output covariance, define
\begin{equation}
 \beta_k
 =
 \frac{v_k}{v_k+\tau^2}.
\end{equation}
Then the output update reduces componentwise to
\begin{align}
 \boxed{
 \mu_{Z_k^{(\Olayer)}\mid\bm y}
 =
 \mu_{Z_k^{(\Olayer)}}
 +
 \beta_k
 \left(
 \mu_{A_k\mid\bm y}-\mu_{Z_k^{(\Olayer)}}
 \right),}
 \label{eq:diagonal-output-mean}
 \\
 \boxed{
 (\sigma_{Z_k^{(\Olayer)}\mid\bm y})^2
 =
 v_k
 +
 \beta_k^2
 \left[
 (\sigma_{A_k\mid\bm y})^2-(v_k+\tau^2)
 \right].}
 \label{eq:diagonal-output-variance}
\end{align}

The exact event posterior is non-Gaussian.  Event AGCI forms the
moment-matched Gaussian approximation
\begin{equation}
 f(\bm z^{(\Olayer)}\mid\bm y)
 \approx
 \Normal\!\left(
 \bm z^{(\Olayer)};
 \bm\mu_{\bm Z^{(\Olayer)}\mid\bm y},
 \bm\Sigma_{\bm Z^{(\Olayer)}\mid\bm y}
 \right).
 \label{eq:moment-projection}
\end{equation}
Standard diagonal TAGI subsequently retains only
$\diag(\bm\Sigma_{\bm Z^{(\Olayer)}\mid\bm y})$.

\section{TAGI backward inference}

Once Event AGCI has supplied the posterior output moments in
\eqref{eq:event-output-mean}--\eqref{eq:event-output-covariance}, the ordinary
TAGI backward pass is unchanged.  Let
\begin{equation}
 \bm Z\equiv\bm Z^{(j)},
 \qquad
 \bm Z^+\equiv\bm Z^{(j+1)}.
\end{equation}
The hidden-unit posterior moments are
\begin{align}
 \bm\mu_{\bm Z\mid\bm y}
 &=
 \bm\mu_{\bm Z}
 +
 \mathbf J_{\bm Z}
 \left(
 \bm\mu_{\bm Z^+\mid\bm y}-\bm\mu_{\bm Z^+}
 \right),
 \\
 \bm\Sigma_{\bm Z\mid\bm y}
 &=
 \bm\Sigma_{\bm Z}
 +
 \mathbf J_{\bm Z}
 \left(
 \bm\Sigma_{\bm Z^+\mid\bm y}-\bm\Sigma_{\bm Z^+}
 \right)
 \mathbf J_{\bm Z}^\trans,
 \\
 \mathbf J_{\bm Z}
 &=
 \bm\Sigma_{\bm Z\bm Z^+}
 \bm\Sigma_{\bm Z^+}^{-1}.
 \label{eq:event-rts-hidden}
\end{align}
For the parameters $\bm\theta\equiv\bm\theta^{(j)}$,
\begin{align}
 \bm\mu_{\bm\theta\mid\bm y}
 &=
 \bm\mu_{\bm\theta}
 +
 \mathbf J_{\bm\theta}
 \left(
 \bm\mu_{\bm Z^+\mid\bm y}-\bm\mu_{\bm Z^+}
 \right),
 \\
 \bm\Sigma_{\bm\theta\mid\bm y}
 &=
 \bm\Sigma_{\bm\theta}
 +
 \mathbf J_{\bm\theta}
 \left(
 \bm\Sigma_{\bm Z^+\mid\bm y}-\bm\Sigma_{\bm Z^+}
 \right)
 \mathbf J_{\bm\theta}^\trans,
 \\
 \mathbf J_{\bm\theta}
 &=
 \bm\Sigma_{\bm\theta\bm Z^+}
 \bm\Sigma_{\bm Z^+}^{-1}.
 \label{eq:event-rts-parameters}
\end{align}

\subsection{Explicit last-layer update}

To avoid confusing the noisy utility $\bm A$ with the activation entering the
last layer, denote the deterministic final-layer feature by
$H_d^{(\Olayer-1)}=h_d$.  The
output unit is
\begin{equation}
 Z_k^{(\Olayer)}
 =
 \sum_d W_{k,d}^{(\Olayer-1)}h_d
 +B_k^{(\Olayer-1)}.
\end{equation}
Under diagonal parameter covariance,
\begin{align}
 \Cov\!\left(
 W_{k,d}^{(\Olayer-1)},Z_k^{(\Olayer)}
 \right)
 &=
 (\sigma_{W_{k,d}^{(\Olayer-1)}})^2h_d,
 \\
 \Cov\!\left(
 B_k^{(\Olayer-1)},Z_k^{(\Olayer)}
 \right)
 &=
 (\sigma_{B_k^{(\Olayer-1)}})^2.
\end{align}
Define the standardized output-moment changes
\begin{align}
 \delta_{\mu,k}
 &=
 \frac{
 \mu_{Z_k^{(\Olayer)}\mid\bm y}
 -\mu_{Z_k^{(\Olayer)}}
 }{v_k},
 \\
 \delta_{v,k}
 &=
 \frac{
 (\sigma_{Z_k^{(\Olayer)}\mid\bm y})^2-v_k
 }{v_k^2}.
\end{align}
The diagonal parameter updates are
\begin{align}
 \mu_{W_{k,d}^{(\Olayer-1)}\mid\bm y}
 &=
 \mu_{W_{k,d}^{(\Olayer-1)}}
 +(\sigma_{W_{k,d}^{(\Olayer-1)}})^2h_d\,\delta_{\mu,k},
 \\
 (\sigma_{W_{k,d}^{(\Olayer-1)}\mid\bm y})^2
 &=
 (\sigma_{W_{k,d}^{(\Olayer-1)}})^2
 +(\sigma_{W_{k,d}^{(\Olayer-1)}})^4h_d^2\,\delta_{v,k},
 \\
 \mu_{B_k^{(\Olayer-1)}\mid\bm y}
 &=
 \mu_{B_k^{(\Olayer-1)}}
 +(\sigma_{B_k^{(\Olayer-1)}})^2\delta_{\mu,k},
 \\
 (\sigma_{B_k^{(\Olayer-1)}\mid\bm y})^2
 &=
 (\sigma_{B_k^{(\Olayer-1)}})^2
 +(\sigma_{B_k^{(\Olayer-1)}})^4\delta_{v,k}.
\end{align}

\section{Mini-batch Event-AGCI}
\label{sec:minibatches}

\subsection{Sequential ADF reference update}

The direct Bayesian interpretation is sequential.  Starting from a prior
$q_0(\bm\theta)$, observation $n$ is assimilated as
\begin{equation}
 q_n(\bm\theta)
 =
 \operatorname{Project}_{\mathcal G}
 \left[
 p(\bm y_n\mid\bm\theta)q_{n-1}(\bm\theta)
 \right],
 \label{eq:sequential-adf}
\end{equation}
where $\operatorname{Project}_{\mathcal G}$ denotes moment projection into the
chosen diagonal Gaussian family.  Each observation is therefore evaluated
using the posterior state produced by the preceding observation.  Batch size
one most closely implements Equation \eqref{eq:sequential-adf}.

\subsection{Shared-state parallel batch approximation}

Let a mini-batch contain $N_b$ deterministic feature vectors, arranged as
\begin{equation}
 \mathbf H_b
 =
 \begin{bmatrix}
 \bm h_1^\trans\\
 \vdots\\
 \bm h_{N_b}^\trans
 \end{bmatrix}
 \in\mathbb R^{N_b\times D}.
\end{equation}
For every example, compute the TAGI output prior and the Event-AGCI posterior
moments using the same pre-batch parameter state.  Collect the standardized
moment changes in
\begin{equation}
 \bm\Delta_\mu
 =
 [\delta_{\mu,nk}]
 \in\mathbb R^{N_b\times K},
 \qquad
 \bm\Delta_v
 =
 [\delta_{v,nk}]
 \in\mathbb R^{N_b\times K}.
\end{equation}
Let $\bm M_{\bm W},\bm S_{\bm W}\in\mathbb R^{K\times D}$ denote the
last-layer weight means and marginal variances, and let
$\bm m_{\bm B},\bm s_{\bm B}\in\mathbb R^K$ denote the bias means and
variances.  Summing the one-example RTS contributions gives the parallel batch
update
\begin{align}
 \boxed{
 \bm M_{\bm W}^{+}
 =
 \bm M_{\bm W}
 +
 \bm S_{\bm W}\odot
 \left(\bm\Delta_\mu^\trans\mathbf H_b\right),}
 \label{eq:batch-weight-mean}
 \\
 \boxed{
 \bm S_{\bm W}^{+}
 =
 \bm S_{\bm W}
 +
 \bm S_{\bm W}^{\odot 2}\odot
 \left[
 \bm\Delta_v^\trans
 (\mathbf H_b\odot\mathbf H_b)
 \right],}
 \label{eq:batch-weight-variance}
 \\
 \boxed{
 \bm m_{\bm B}^{+}
 =
 \bm m_{\bm B}
 +
 \bm s_{\bm B}\odot
 \left(\bm\Delta_\mu^\trans\bm 1_{N_b}\right),}
 \label{eq:batch-bias-mean}
 \\
 \boxed{
 \bm s_{\bm B}^{+}
 =
 \bm s_{\bm B}
 +
 \bm s_{\bm B}^{\odot 2}\odot
 \left(\bm\Delta_v^\trans\bm 1_{N_b}\right).}
 \label{eq:batch-bias-variance}
\end{align}
Here $\odot$ denotes elementwise multiplication and
$\bm S^{\odot 2}=\bm S\odot\bm S$.

For ordinary one-pass evidence assimilation, the per-example contributions in
Equations \eqref{eq:batch-weight-mean}--\eqref{eq:batch-bias-variance} are
\emph{summed}, not averaged.  Replacing the sum by a batch average weakens each
batch likelihood by approximately a factor $1/N_b$ and is closer to a tempered
or generalized-Bayes update.

\subsection{What batching approximates}

The parallel update is not equal to sequential ADF because all examples use a
shared pre-batch state.  In particular, it neglects:
\begin{enumerate}
 \item the change in parameter moments after each observation;
 \item posterior correlations between the outputs of different examples that
 arise because they share the same random parameters;
 \item cross-class and cross-parameter covariance removed by the diagonal
 projection.
\end{enumerate}
Consequently, smaller batches are closer to sequential ADF, whereas larger
batches provide greater computational throughput and average event-specific
covariance fluctuations.

For a fixed input, the Event-AGCI covariance obeys
\begin{equation}
 \E_{Y\sim\operatorname{Categorical}(\bm p)}
 \left[
 \bm\Sigma_{\bm Z^{(\Olayer)}\mid Y}
 \right]
 =
 \sum_{c=1}^{K}p_c
 \bm\Sigma_{\bm Z^{(\Olayer)}\mid Y=c}.
 \label{eq:event-covariance-average}
\end{equation}
The right-hand side is the residual covariance of full LMMSE-AGCI.  A shuffled
batch therefore makes the observed-event covariance update behave like a
stochastic estimate of the class-averaged covariance when the model
probabilities are close to the data-label probabilities.  This explains why
large, representative batches can make Event AGCI and full LMMSE-AGCI nearly
indistinguishable after one pass.

\subsection{Recommended batching protocol}

A defensible protocol for a frozen-feature TAGI last layer is:
\begin{enumerate}
 \item compute and freeze all feature preprocessing using training data only;
 \item initialize the prior once using Section \ref{sec:initialization};
 \item globally shuffle examples before assimilation;
 \item use batches that are approximately representative of the full class
 distribution;
 \item compute every Event-AGCI moment from the same pre-batch state;
 \item sum the per-example updates and commit one posterior state per batch;
 \item verify that all updated variances remain positive;
 \item assimilate every training observation once for the cleanest ordinary
 Bayesian interpretation.
\end{enumerate}
Class-ordered batches should be avoided because ADF is order-dependent.  If a
variance update becomes invalid, first reduce the batch size or prior scale.
A coordinate-wise cap is a numerical heuristic and changes the stated
moment-matching update; if used frequently, it should be reported.

Repeated passes over the same data do not represent repeated refinement of one
ordinary posterior.  Ideally, $E$ passes correspond to the power posterior
\begin{equation}
 p_E(\bm\theta\mid\mathcal D)
 \propto
 p(\bm\theta)
 p(\mathcal D\mid\bm\theta)^E,
\end{equation}
with additional discrepancy caused by ADF projection and batching.  Multiple
passes should therefore be described as tempered or generalized-Bayes
inference rather than ordinary one-pass Bayes.

\section{Gauss--Hermite quadrature}

Let $x_q$ and $w_q$, $q=1,\ldots,Q$, be Gauss--Hermite nodes and weights for
integrals with weight $\exp(-x^2)$.  Set
\begin{equation}
 a_q=\mu_c+\sqrt{2}\,r_c x_q.
\end{equation}
Equation \eqref{eq:class-probability} is approximated by
\begin{equation}
 p_c
 \approx
 \frac{1}{\sqrt\pi}
 \sum_{q=1}^{Q}
 w_q
 \prod_{j\ne c}
 \Phi\!\left(
 \frac{a_q-\mu_j}{r_j}
 \right).
 \label{eq:gh-probability}
\end{equation}
Define normalized event weights
\begin{equation}
 \omega_q
 =
 \frac{
 w_q\prod_{j\ne c}\Phi(\alpha_j(a_q))
 }{
 \sum_{\ell=1}^{Q}
 w_\ell\prod_{j\ne c}\Phi(\alpha_j(a_\ell))
 }.
 \label{eq:gh-event-weights}
\end{equation}
The integrals in Sections 5 and 6 are then evaluated as weighted sums.  For
example,
\begin{align}
 \mu_{A_c\mid\bm y}
 &\approx
 \sum_q\omega_qa_q,
 \\
 \mu_{A_j\mid\bm y}
 &\approx
 \sum_q\omega_qm_j(a_q),
 \\
 \E[A_j^2\mid Y=c]
 &\approx
 \sum_q\omega_q\left[s_j^2(a_q)+m_j(a_q)^2\right].
\end{align}
For numerical stability, the products of Gaussian CDFs should be formed in
the log domain, and $\lambda_j(a)$ should be evaluated using
$\log\phi(\alpha_j)-\log\Phi(\alpha_j)$.

\section{Efficient diagonal Event-AGCI algorithm}

For one labeled example $(\bm x,c)$, the diagonal algorithm is:
\begin{enumerate}
 \item Perform the TAGI forward pass and obtain
 $\bm\mu_{\bm Z^{(\Olayer)}}$ and
 $\diag(\bm\Sigma_{\bm Z^{(\Olayer)}})$.

 \item Compute $r_k^2=v_k+\tau^2$ for every class.

 \item Use the observed winner $c$ as the one-dimensional quadrature variable
 and compute the normalized event weights in
 \eqref{eq:gh-event-weights}.

 \item Compute $\mu_{A_k\mid\bm y}$ and
 $(\sigma_{A_k\mid\bm y})^2$ for every class using
 \eqref{eq:winner-mean}--\eqref{eq:loser-variance}.

 \item Transform those moments to
 $\mu_{Z_k^{(\Olayer)}\mid\bm y}$ and
 $(\sigma_{Z_k^{(\Olayer)}\mid\bm y})^2$ using
 \eqref{eq:diagonal-output-mean}--\eqref{eq:diagonal-output-variance}.

 \item Propagate the output posterior moments backward with the ordinary TAGI
 equations \eqref{eq:event-rts-hidden}--\eqref{eq:event-rts-parameters}.
\end{enumerate}
For $Q$ quadrature nodes, the diagonal observed-event calculation is
$O(QK)$ per example and does not require a dense $K\times K$ event covariance.

\section{Relationship with full LMMSE-AGCI}

Let
\begin{equation}
 \bm p=(p_1,\ldots,p_K)^\trans,
 \qquad
 \bm m_c=\E[\bm Z^{(\Olayer)}\mid Y=c].
\end{equation}
The categorical observation covariance is
\begin{equation}
 \bm\Sigma_{\bm Y}
 =
 \diag(\bm p)-\bm p\bm p^\trans,
\end{equation}
and the output--label cross-covariance has columns
\begin{equation}
 [\bm\Sigma_{\bm Z^{(\Olayer)}\bm Y}]_{:,c}
 =
 p_c
 \left(
 \bm m_c-\bm\mu_{\bm Z^{(\Olayer)}}
 \right).
\end{equation}
The full categorical LMMSE gain is
\begin{equation}
 \mathbf J_{\mathrm{AGCI}}
 =
 \bm\Sigma_{\bm Z^{(\Olayer)}\bm Y}
 \bm\Sigma_{\bm Y}^{\dagger}.
\end{equation}
Its mean update is
\begin{equation}
 \bm\mu_{\bm Z^{(\Olayer)}}
 +
 \mathbf J_{\mathrm{AGCI}}
 (\bm e_c-\bm p)
 =
 \bm m_c.
\end{equation}
Thus Event AGCI and full LMMSE-AGCI have the same class-conditional mean when
their moments are evaluated consistently.  They differ in their covariance:
\begin{align}
 \text{Event AGCI:}\quad
 &\bm\Sigma_{\bm Z^{(\Olayer)}\mid Y=c},
 \\
 \text{Full LMMSE-AGCI:}\quad
 &\sum_{c=1}^{K}
 p_c\bm\Sigma_{\bm Z^{(\Olayer)}\mid Y=c}.
\end{align}
Consequently,
\begin{equation}
 \E_{Y\sim\operatorname{Categorical}(\bm p)}
 \left[
 \bm\Sigma_{\bm Z^{(\Olayer)}\mid Y}
 \right]
 =
 \bm\Sigma_{\bm Z^{(\Olayer)},\mathrm{LMMSE}}^+.
\end{equation}
Event AGCI therefore uses the covariance associated with the observed event,
whereas full LMMSE-AGCI uses the class-averaged residual covariance.

There is no separate squared-label channel in categorical AGCI because a
one-hot observation satisfies $\bm Y\odot\bm Y=\bm Y$.  Its second-moment
information is already contained in
$\diag(\bm p)-\bm p\bm p^\trans$.

\section{Approximation boundary}

\begin{center}
\begin{tabular}{>{\raggedright\arraybackslash}p{0.43\textwidth}
                >{\raggedright\arraybackslash}p{0.43\textwidth}}
\toprule
Exact under the stated model & Approximate in the implemented method \\
\midrule
Gaussian random-utility likelihood
& Numerical Gauss--Hermite integration \\
One-dimensional event representation for diagonal utilities
& Gaussian projection of the non-Gaussian event posterior \\
Gaussian transformation from $\bm A\mid\bm y$ to the first two moments of
$\bm Z^{(\Olayer)}\mid\bm y$
& Removal of event-induced cross-class covariance \\
Class-specific output moments up to quadrature error
& Diagonal parameter covariance and layer-wise conditional independence \\
One-observation TAGI/RTS moment propagation
& Shared-state mini-batches, observation ordering, repeated passes, and any
update cap or damping \\
\bottomrule
\end{tabular}
\end{center}

\section{Implementation checks}

The following checks should hold up to numerical quadrature error:
\begin{enumerate}
 \item Predictive probabilities form a partition:
 \begin{equation}
  \sum_{c=1}^{K}p_c=1.
 \end{equation}

 \item Posterior covariance is symmetric and positive semidefinite:
 \begin{equation}
  \bm\Sigma_{\bm Z^{(\Olayer)}\mid\bm y}
  =
  \bm\Sigma_{\bm Z^{(\Olayer)}\mid\bm y}^{\trans},
  \qquad
  \lambda_{\min}
  (\bm\Sigma_{\bm Z^{(\Olayer)}\mid\bm y})
  \ge -\varepsilon.
 \end{equation}

 \item Translation invariance holds for any scalar $a$:
 \begin{equation}
  p_c(\bm\mu+a\bm 1,\bm\Sigma,\tau)
  =p_c(\bm\mu,\bm\Sigma,\tau).
 \end{equation}

 \item Joint scale invariance holds for $a>0$:
 \begin{equation}
  p_c(a\bm\mu,a^2\bm\Sigma,a\tau)
  =p_c(\bm\mu,\bm\Sigma,\tau).
 \end{equation}

 \item For $K=2$, the quadrature moments agree with the analytical moments of
 a univariate truncated Gaussian utility difference.
\end{enumerate}

\section{Summary}

Event AGCI replaces the linear Gaussian TAGI observation model with a proper
categorical likelihood induced by Gaussian random utilities.  For the observed
class, it computes the first two moments of the class-winning event, transforms
those moments into posterior moments for the TAGI output units, and propagates
them backward using the unchanged TAGI RTS equations.  In compact form,
\begin{equation}
\boxed{
\bm Z^{(\Olayer)}
\xrightarrow{\;\bm A=\bm Z^{(\Olayer)}+\bm E\;}
Y=\arg\max_k A_k
\xrightarrow{\;\text{event moments}\;}
f(\bm z^{(\Olayer)}\mid\bm y)
\xrightarrow{\;\text{TAGI RTS}\;}
f(\bm\theta\mid\bm y).}
\end{equation}

\end{document}
